
\subsection{Related Work}\label{sec:related-work}

Goodhart’s law  was first introduced by \citet{goodhart_problems_1975}, and has later been elaborated upon by works such as \citet{manheim_categorizing_2019}.
Goodhart's law has also previously been studied in the context of machine learning. 
In particular, \citet{hennessy_goodharts_2023} investigate Goodhart's law analytically in the context where a machine learning model is used to evaluate an agent’s actions -- unlike them, we specifically consider the RL setting\oliver{I'm guessing what we're doing is different in other, more useful ways as well?}. \citet{ashton_causal_2021} shows by example that RL systems can be susceptible to Goodharting in certain situations. In contrast, we show that Goodhart's law is a robust phenomenon across a wide range of environments, explain why it occurs in RL, and use it to devise new solution methods.

In the context of RL, Goodhart's law is closely related to \emph{reward gaming}. Specifically, if reward gaming means an agent finding an unintended way to increase its reward, then Goodharting is an \emph{instance} of reward gaming where optimisation of the proxy initially leads to desirable behaviour, followed by a decrease after some threshold.
\Citet{krakovna_specification_2020} list illustrative examples of reward hacking, while \citet{reward_misspecification} manually construct proxy rewards for several environments and then demonstrate that most of them lead to reward hacking.
\citet{subset_features} consider proxy rewards that depend on a strict subset of the features which are relevant to the true reward and then show that optimising such a proxy in some cases may be arbitrarily bad, given certain assumptions.
\cite{skalse_defining_2022} introduce a theoretical framework for analysing reward hacking. They then demonstrate that, in any environment and for any true reward function, it is impossible to create a non-trivial proxy reward that is guaranteed to be unhackable. 
% (unless the proxy is equivalent to the true reward)
Also relevant, \citet{reward_corruption} study the related problem of reward corruption, \citet{song_observational_2019} investigate overfitting in model-free RL due to faulty implications from correlations in the environment, and
\citet{pang_reward_2023} examine reward gaming in language models. 
Unlike these works, we analyse reward hacking through the lens of \emph{Goodhart's law} and show that this perspective provides novel insights.

\cite{gao_scaling_2022} consider the setting where a large language model is optimised against a reward model that has been trained on a \enquote{gold standard} reward function, and investigate how the performance of the language model according to the gold standard reward scales in the size of the language model, the amount of training data, and the size of the reward model.
They find that the performance of the policy follows a Goodhart curve, where the slope gets less prominent %and the performance gets better 
for larger reward models and larger amounts of training data.
Unlike them, we do not only focus on language, but rather, aim to establish to what extent Goodhart dynamics occur for a wide range of RL environments. Moreover, we also aim to \emph{explain} Goodhart's law, and use it as a starting point for developing new algorithms.
